Questa tesi ha proposto e valutato RowMask23, una parametrizzazione strutturata dei pesi per le proiezioni lineari dei modelli linguistici Transformer, e l'ha applicata a un modello addestrato da zero, STERA 30M. In questo capitolo riassumiamo l'esito delle quattro ipotesi, discutiamo i limiti del lavoro e indichiamo le direzioni in cui proseguirlo.

\section{Esito delle ipotesi}

\begin{table}[htbp]
\centering
\small
\begin{tabularx}{\linewidth}{lX>{\raggedright\arraybackslash}p{3.4cm}}
\toprule
& Evidenza principale & Esito \\
\midrule
H1 & NLL TEST 1,174, inferiore a TinyStories-8M (1,515) e TinyStories-33M (1,324) & Confermata \\
H2 & 121,385~MB in FP32 contro 330,885~MB dell'espansione densa ($-63{,}3\%$) e 274,056~MB di TinyStories-33M ($-55{,}7\%$) & Confermata \\
H3 & Decode $-36{,}1\%$ rispetto al motore denso sugli stessi pesi, con token identici & Confermata \\
H4 & 44,0\% di racconti corretti contro 21,7\% e 48,3\% & Confermata rispetto a TinyStories-8M; non conclusiva rispetto a TinyStories-33M \\
\bottomrule
\end{tabularx}
\caption{Sintesi delle ipotesi della tesi e del loro esito.}
\label{tab:esito-ipotesi}
\end{table}

\paragraph{Una famiglia di matrici con buone proprietà.} Un blocco RowMask23 da $128\times128$ usa \num{3176} coordinate anziché \num{16384} coefficienti indipendenti. La mappa dalle coordinate alla matrice è lineare e iniettiva, il gradiente rispetto alle coordinate ha una forma chiusa e il prodotto matrice--vettore si calcola con la trasformata di Walsh--Hadamard senza materializzare la matrice.

\paragraph{Un modello che apprende.} STERA 30M, addestrato da zero con tutte le proiezioni vincolate a RowMask23, raggiunge una NLL di 1,152 sul campione DEV e di 1,174 sui racconti TEST. La curva di addestramento non mostra segni di sovradattamento.

\paragraph{Racconti fra i due riferimenti.} Nel giudizio cieco, STERA 30M scrive racconti corretti nel 44\% dei casi: il doppio di TinyStories-8M, con cui condivide l'ordine di grandezza dei parametri memorizzati, e 4 punti in meno di TinyStories-33M, che ne ha più del doppio, con una differenza non significativa.

\paragraph{Un costo minore in memoria.} Il risparmio di byte è una conseguenza diretta del conteggio dei parametri ed è stato verificato sui pesi effettivamente salvati. Sullo stesso modello, il motore compatto riduce il tempo per token in decode, mentre il motore denso resta più rapido nell'elaborazione del prompt.

\section{Limiti}

\paragraph{Un solo addestramento.} STERA 30M è stato addestrato una volta, con un unico seme. Gli intervalli di confidenza riportati misurano l'incertezza dovuta ai racconti di prova, non quella dovuta all'inizializzazione.

\paragraph{Riferimenti non controllati.} I checkpoint TinyStories sono modelli pubblici con architettura, budget di addestramento e dati non identici a quelli di STERA 30M. Il confronto descrive le prestazioni dei modelli esistenti e non isola l'effetto della sola parametrizzazione; in particolare, il vantaggio di NLL può dipendere in parte dalla corrispondenza fra i dati di addestramento di STERA 30M e i racconti di prova.

\paragraph{Calcolo e memoria.} Il vantaggio di RowMask23 riguarda i byte da memorizzare e trasferire, non il numero di operazioni: durante l'addestramento e nel motore denso, STERA 30M calcola come un modello con \num{82721280} valori.

\paragraph{Valutazione automatica.} La qualità dei racconti è stata giudicata da un singolo modello linguistico, in modo cieco e con una rubrica fissata in anticipo. Una validazione con lettori umani renderebbe la stima più robusta.

\paragraph{Misure di sistema.} Le misure di esecuzione sono state effettuate su un solo dispositivo, un Mac con chip M1 e 16~GiB di memoria unificata.

\section{Sviluppi futuri}

\paragraph{Un confronto controllato.} Il passo più importante è addestrare, con gli stessi dati, lo stesso budget e più semi, STERA 30M e due modelli densi di riferimento: uno con lo stesso numero di parametri memorizzati e uno con la stessa forma. Questo confronto isolerebbe l'effetto della parametrizzazione.

\paragraph{Esecuzione fuori memoria.} Una misura completa di generazione con lo stesso modello eseguito in forma compatta residente e in forma densa trasferita dal disco quantificherebbe direttamente il vantaggio di RowMask23 nello scenario che motiva questo lavoro.

\paragraph{Kernel più efficienti.} Memorizzare la trasformata della base, fondere in un unico kernel le proiezioni di uno strato e scrivere un kernel diretto anche per l'elaborazione del prompt sono ottimizzazioni immediate.

\paragraph{Struttura e quantizzazione.} La struttura riduce il numero di valori, la quantizzazione il numero di bit per valore. Un addestramento consapevole della quantizzazione delle coordinate permetterebbe di combinare i due effetti.

\paragraph{Altre famiglie della stessa specie.} Il numero di correzioni per riga e la regola che ne fissa le posizioni sono scelte di progetto. Famiglie con più o meno correzioni offrirebbero un continuo di compromessi fra compressione e capacità, da esplorare con la stessa metodologia.

\section{Considerazioni finali}

I modelli linguistici sono sempre più eseguiti su dispositivi personali, dove la memoria resta la risorsa più scarsa. Questa tesi mostra che si può affrontare il problema scegliendo la forma dei pesi prima dell'addestramento, invece di comprimerli dopo: un Transformer con proiezioni RowMask23 apprende, genera racconti la cui qualità si colloca fra quella di due modelli densi di riferimento, senza differenze significative rispetto al più grande, e occupa una frazione della sua memoria.
